% =============================================================================
%                    WEEK 1: INTRODUCTION
% =============================================================================
\section{Week 1: Introduction}

% -----------------------------------------------------------------------------
%                              1.1 TOPIC SUMMARY
% -----------------------------------------------------------------------------
\subsection{Topic Summary}

\begin{keybox}[Learning Objectives]
By the end of this week, you will be able to:
\begin{enumerate}
  \item Understand what \textbf{philosophy} is and its major branches.
  \item Explain why philosophy is \textbf{relevant to AI}.
  \item Describe how philosophical issues impact the \textbf{ethics of AI}.
  \item Understand the \textbf{structure and assessment} of this module.
\end{enumerate}
\end{keybox}

\separator

\subsubsection*{The Big Picture}

This introductory lecture establishes why philosophical thinking matters for AI ethics. Philosophy provides the conceptual tools---about mind, knowledge, ethics, and logic---that we need to address questions like: Can machines think? Should autonomous systems be held responsible? How do we ensure AI benefits humanity?

The module is divided into two halves: (1) philosophical foundations (intelligence, consciousness, reasoning, ethics), and (2) applied AI ethics (algorithms, medicine/war, superintelligence, society).

\separator

\subsubsection*{What Examiners Like to Ask}

\begin{itemize}
  \item Define the major branches of philosophy (metaphysics, epistemology, ethics, logic).
  \item Explain Descartes' argument for mind-body dualism and its implications.
  \item How does determinism challenge traditional notions of moral responsibility?
  \item Demonstrate how logic can be used to analyse ethical arguments (Socratic method).
  \item Why are these philosophical concepts relevant to AI ethics?
\end{itemize}

\bigseparator

% -----------------------------------------------------------------------------
%                 1.2 OVERVIEW + CORE CONCEPTS + EXTENDED THEORY
% -----------------------------------------------------------------------------
\subsection{Overview + Core Concepts + Extended Theory}

\subsubsection*{Roadmap}
\begin{enumerate}
  \item What is philosophy?
  \item Major branches of philosophy
  \item Example 1: Descartes and mind-body dualism
  \item Example 2: Free will, determinism, and responsibility
  \item Example 3: Logic and ethical argument (Socratic method)
  \item Relevance to AI ethics
  \item Module structure overview
\end{enumerate}

\bigseparator

% --- What is Philosophy? ---
\subsubsection{What is Philosophy?}

\begin{defbox}[Philosophy]
\textbf{Philosophy} (literally ``love of wisdom'') is an activity people undertake when they seek to understand \textbf{fundamental truths} about:
\begin{itemize}
  \item Themselves (What is the self? What is consciousness?)
  \item The world (What exists? What is real?)
  \item Their relationships to the world and each other (What is right? How should we live?)
\end{itemize}
\end{defbox}

Philosophy addresses life's most basic questions:
\begin{itemize}
  \item Why are we here? Why is there something rather than nothing?
  \item What is knowledge? How do we know what we know?
  \item What is consciousness? What is the self?
  \item What is the good life? What are morally right actions?
\end{itemize}

\bigseparator

% --- Branches of Philosophy ---
\subsubsection{Major Branches of Philosophy}

\begin{center}
\renewcommand{\arraystretch}{1.4}
\begin{tabular}{@{} l p{5cm} p{5cm} @{}}
\toprule
\textbf{Branch} & \textbf{Focus} & \textbf{Key Questions} \\
\midrule
\textbf{Metaphysics} & Nature of reality, existence, being & Is there a God? What is truth? What is a person? Is the world composed of matter only? Do people have minds? \\
\textbf{Epistemology} & Nature of knowledge & What is knowledge? How do we know anything? How can we justify beliefs? \\
\textbf{Ethics} & What we ought to do & What is good? What makes actions right? Is morality objective or subjective? \\
\textbf{Logic} & Structure of arguments & What constitutes good reasoning? How do we evaluate arguments? \\
\bottomrule
\end{tabular}
\end{center}

\begin{tipbox}[Exam Relevance]
These four branches recur throughout the module:
\begin{itemize}
  \item \textbf{Metaphysics} $\to$ Can machines have minds? (Weeks 2--3)
  \item \textbf{Epistemology} $\to$ Can machines know things? (Week 4)
  \item \textbf{Ethics} $\to$ What should AI do? (Weeks 5--9)
  \item \textbf{Logic} $\to$ How do we reason about ethics? (Week 4)
\end{itemize}
\end{tipbox}

\bigseparator

% --- Example 1: Descartes ---
\subsubsection{Example 1: Descartes and Mind-Body Dualism}

\begin{thmbox}[Descartes' Cogito Argument (1637)]
\textbf{``I think, therefore I am''} (\textit{Cogito ergo sum})

\textbf{The argument:}
\begin{enumerate}
  \item The very act of \textbf{thinking} proves the existence of the \textbf{self} (the ``I'' doing the thinking).
  \item Everything else---the external world, our bodies---could be an illusion (``the Matrix'').
  \item We can be \textbf{absolutely certain} of the self, but \textbf{not certain} of the physical world.
  \item Therefore, the \textbf{mind} (certain) must be fundamentally \textbf{different} from the \textbf{body} (uncertain).
\end{enumerate}

\textbf{Conclusion:} Mind and body are separate substances (\textbf{dualism}).
\end{thmbox}

\begin{defbox}[Mind-Body Dualism]
The view that \textbf{mind} and \textbf{body} are fundamentally different kinds of things:
\begin{itemize}
  \item Mind: non-physical, mental, private, certain
  \item Body: physical, material, public, uncertain
\end{itemize}
\end{defbox}

\begin{exbox}[Impact on Western Medicine]
Cartesian dualism dominated Western thought for centuries:
\begin{itemize}
  \item Traditional view: Mind and body are \textbf{separate}, so mental processes \textbf{cannot affect} physical health.
  \item ``Positive thinking'' affecting health was \textbf{dismissed}.
  \item Modern view: Mind and body are \textbf{interconnected}. Mental states (stress, optimism) \textbf{do affect} physical health.
  \item Eastern medicine traditions (Buddhist, Ayurvedic) always recognised this connection.
\end{itemize}
\end{exbox}

\begin{tipbox}[Relevance to AI]
The mind-body problem is central to AI:
\begin{itemize}
  \item Can a \textbf{physical machine} have a \textbf{mind}?
  \item If mind and body are separate, could machines ever be conscious?
  \item If mind emerges from physical processes, perhaps machines \textit{can} think.
\end{itemize}
We return to this in Weeks 2 (Intelligence) and 3 (Consciousness).
\end{tipbox}

\bigseparator

% --- Example 2: Free Will ---
\subsubsection{Example 2: Free Will, Determinism, and Responsibility}

\begin{defbox}[Determinism]
\textbf{Determinism:} All events, including human decisions, are determined by prior causes according to physical laws.

\textbf{Implication:} If physics is deterministic (A causes B causes C\ldots), and mental processes are ultimately physical (neural activity), then:
\begin{itemize}
  \item Every decision we make is the \textbf{effect} of a chain of causes.
  \item \textbf{Free will is an illusion}---there is no ``self'' that freely chooses independently of physics.
\end{itemize}
\end{defbox}

\begin{thmbox}[Implications for Justice]
If there is no free will, then when a criminal commits a crime:
\begin{itemize}
  \item There is \textbf{no self} that ``freely chose'' to commit the crime.
  \item The crime was a \textbf{deterministic consequence} of prior causes (genetics, environment, neural processes).
  \item \textbf{Punishment} (retribution against a morally responsible self) may be unjustified.
\end{itemize}

\textbf{Alternative justifications for penalties:}
\begin{enumerate}
  \item \textbf{Rehabilitation:} Change the criminal's decision-making so future decisions differ.
  \item \textbf{Deterrence:} Signal to others that crime leads to penalty.
  \item \textbf{Protection:} Protect society from the criminal (and vice versa).
\end{enumerate}
\end{thmbox}

\begin{tipbox}[Relevance to AI]
These questions apply directly to autonomous AI:
\begin{itemize}
  \item When an \textbf{autonomous car} kills a pedestrian, who is responsible?
  \item When a \textbf{lethal autonomous weapon} kills, who is accountable?
  \item Should we ``punish'' the AI? The designers? The users?
  \item How do we design penalties/incentives to prevent future harms?
\end{itemize}
We return to this in Week 6 (Algorithms) and Week 7 (AI in War).
\end{tipbox}

\bigseparator

% --- Example 3: Logic and Ethics ---
\subsubsection{Example 3: Logic and Ethical Argument}

\begin{defbox}[The Socratic Method]
\philosopher{Socrates} used a method of \textbf{questioning} to expose \textbf{contradictions} in his interlocutors' arguments.

\textbf{Key idea:} By revealing inconsistencies, we can refine our beliefs and arrive at better-justified positions.
\end{defbox}

\begin{exbox}[Socrates vs.\ Sara: The Abortion Debate]
\textbf{Dialogue:}
\begin{enumerate}
  \item Sara: ``Abortion is wrong.''
  \item Socrates: ``Why?''
  \item Sara: ``Human life is sacred, and a fetus is human.''
  \item Socrates: ``So?''
  \item Sara: ``If a life is sacred, it should not be terminated.''
  \item Socrates: ``So you should be against the death penalty for murderers?''
  \item Sara: ``Err\ldots no! I'm in favour.''
\end{enumerate}

\textbf{Sara's implicit reasoning:}
\begin{enumerate}
  \item If X is human $\to$ X's life is sacred.
  \item A fetus is human $\to$ A fetus's life is sacred.
  \item If X's life is sacred $\to$ X should not be terminated.
  \item A fetus's life is sacred $\to$ A fetus should not be terminated.
\end{enumerate}

\textbf{The contradiction:} The same reasoning applies to murderers (who are also human), so Sara should oppose the death penalty---but she supports it!
\end{exbox}

\begin{thmbox}[Refining Knowledge Through Refutation]
\textbf{Solution:} Sara can refine her rule to resolve the contradiction:

\textbf{Old rule:} If X is human $\to$ X's life is sacred.

\textbf{Refined rule:} If X is human \textbf{and} X does not take the life of another human $\to$ X's life is sacred.

Now Sara can consistently:
\begin{itemize}
  \item Oppose abortion (fetus is human and hasn't killed anyone).
  \item Support the death penalty (murderer has taken a life).
\end{itemize}

\textbf{Key insight:} Knowledge evolves through refutation---we refine rules to accommodate counterexamples.
\end{thmbox}

\begin{tipbox}[Relevance to AI]
\begin{itemize}
  \item \textbf{AI systems may need to reason about ethics} using logic.
  \item \textbf{Argumentation} (exchange of argument and counter-argument) is a formal method for ethical reasoning.
  \item We may use \textbf{AI-assisted reasoning} to help humans make better moral decisions.
\end{itemize}
We return to this in Week 4 (Reasoning) and Week 5 (Ethics).
\end{tipbox}

\bigseparator

% --- Module Structure ---
\subsubsection{Module Structure Overview}

\begin{center}
\renewcommand{\arraystretch}{1.3}
\begin{tabular}{@{} c l p{6.5cm} @{}}
\toprule
\textbf{Week} & \textbf{Topic} & \textbf{Key Questions} \\
\midrule
1 & Introduction & What is philosophy? Why relevant to AI? \\
2 & Intelligence & What is intelligence? Can machines be intelligent? \\
3 & Consciousness & What is consciousness? Could AI be conscious? \\
4 & Reasoning \& Communication & Symbolic vs.\ ML paradigms; argumentation \\
5 & Ethics \& Morality & Ethical theories; implementing moral agents \\
6 & Algorithms & Bias, accountability, responsibility \\
7 & AI in Medicine \& War & Lethal autonomous weapons; medical AI ethics \\
8 & Superintelligence & AGI $\to$ superintelligence; value loading problem \\
9 & AI \& Human Society & Echo chambers; surveillance capitalism; moral progress \\
\bottomrule
\end{tabular}
\end{center}

\bigseparator

% --- Summary ---
\subsubsection{Summary: Key Takeaways}

\begin{thmbox}[Week 1 Key Points]
\begin{enumerate}
  \item \textbf{Philosophy} addresses fundamental questions about reality, knowledge, ethics, and reasoning.
  \item \textbf{Descartes' dualism} separated mind from body---this matters for whether machines can have minds.
  \item \textbf{Determinism} challenges free will---this matters for assigning responsibility to AI systems.
  \item \textbf{Logic and the Socratic method} help us analyse and refine ethical arguments.
  \item These philosophical tools will be applied throughout the module to AI ethics questions.
\end{enumerate}
\end{thmbox}

\begin{defbox}[Key Terms to Remember]
\begin{itemize}
  \item \textbf{Metaphysics:} Study of the nature of reality and existence.
  \item \textbf{Epistemology:} Study of knowledge and how we know things.
  \item \textbf{Ethics:} Study of what we ought to do; right and wrong.
  \item \textbf{Logic:} Study of the structure and validity of arguments.
  \item \textbf{Dualism:} Mind and body are fundamentally different substances.
  \item \textbf{Determinism:} All events are caused by prior events according to physical laws.
  \item \textbf{Modus ponens:} If P then Q; P; therefore Q.
  \item \textbf{Refutation:} Revising beliefs when contradictions are discovered.
\end{itemize}
\end{defbox}

